% META: {"id":"1SPE-GEOREP-EX-043","chapitre":"1SPE-GEOMETRIE-REPEREE","type_objet":"exercice","capacites":["1SPE-GEOMETRIE-REPEREE-C5"],"capacites_codes":["C5"],"methodes":["M5"],"parcours":2,"competences":["calculer","raisonner"],"duree_min":15,"mode_creation":"ex_nihilo","sources_inspiration":[],"fichier_tex":"chapitres/1SPE-GEOMETRIE-REPEREE/exercices/1SPE-GEOREP-EX-043.tex","corrige_tex":"chapitres/1SPE-GEOMETRIE-REPEREE/corriges/1SPE-GEOREP-CO-043.tex","status":"approved","gestes":["raisonnement_par_coordonnees"],"origin":"generated","status_history":["generated","audited","accepted","approved"],"release_acceptance":"RELEASE_OWNER_BATCH_ACCEPTANCE","acceptance_closure_digest":"sha256:9b3ccf9a81c5520fbb7e03b7b2d3e7bf2057a02834b6d908bf3be5f49f3b553f","acceptance_receipt_digest":"sha256:4fad86f0282e0fa4e0419cca1ad6379ae7af5a280c04a4a90880f90a1c044242"}
% BEGIN-VERIFY
% from sympy import *
% x, y = symbols('x y')
% # Projete orthogonal de P(4,5) sur D: 3x + 4y - 6 = 0
% # Perpendiculaire: 4x - 3y + c = 0 passant par P: 16-15+c=0 => c=-1
% # Systeme: 3x+4y=6, 4x-3y=1
% sol = solve([3*x+4*y-6, 4*x-3*y-1], [x,y])
% assert sol[x] == Rational(22,25)
% assert sol[y] == Rational(21,25)
% END-VERIFY
\begin{exercice}{1SPE-GEOREP-EX-043}{2}{15}
Déterminer le projeté orthogonal du point $P(4 ; 5)$ sur la droite $\mathcal{D} : 3x + 4y - 6 = 0$.
\end{exercice}
